\subsubsection{替换函数检验与修正模块}

替换函数检验与修正模块是系统的质量守门员，负责确保替换代码的正确性和可靠性。该模块基于 VerifyReplacement 工具实现，由两部分检验机制共同构成：代码审查规则验证和编译验证（可选）。这两部分机制通过不同的验证维度，从代码规则和编译可执行性两个层面共同保障替换代码的质量。

\textbf{代码审查规则验证}是替换代码质量的第一道防线，通过预定义的规则集检查代码的语法正确性和语义一致性。该验证机制包括四条核心规则：签名一致性规则要求替换函数的名称、返回类型、参数列表必须与原函数完全一致，这是替换代码能够正确集成回工程的基本前提；语义保持性规则要求替换代码应保留原函数的核心业务逻辑，如状态变量更新、错误处理流程等，只替换与硬件交互的部分；代码合理性规则要求替换代码不应引入明显的问题，如未声明变量、无限循环、内存泄漏等；调用关系保持规则要求如果原函数调用了某些核心函数，替换代码应在适当的位置保留这些调用。代码审查规则验证采用 fail closed 策略，当验证失败时直接拒绝持久化替换代码，确保替换代码的质量底线，避免错误代码进入后续流程。

\textbf{编译验证}（可选）是替换代码质量的第二道防线，通过实际编译测试验证替换代码的可执行性。编译验证机制采用临时应用和最终恢复策略，在不污染工作树的前提下验证替换代码的编译兼容性。临时应用是将替换代码临时应用到源文件，替换原函数的定义，确保验证环境与实际使用环境的一致性。构建验证是在临时应用的状态下，调用构建系统执行完整的编译链接流程，检测替换代码的语法错误、类型错误、链接错误等各类编译错误。最终恢复是不论编译结果如何，都要恢复原始函数定义，确保工作树不被污染。编译验证失败时，VerifyReplacement 工具会返回编译错误信息，系统可以触发自动修复机制，或者直接拒绝该替换代码，确保只有经过严格验证的代码才能被采纳。
